\section{Rief\/fel's pseudodif\/ferential calculus: a short review}\label{sectra}

We start by describing brief\/ly Rief\/fel deformation \cite{Rie1,Rie2} in a slightly restricted setting.
The initial object, containing {\it the classical data}, is a quadruplet
$\left(\A,\Theta,\Xi,[\![\cdot,\cdot,]\!]\right)$. The pair $\left(\Xi,[\![\cdot,\cdot]\!]\right)$ will
be taken to be a $2n$-dimensional symplectic vector space.
On the other hand $\left(\A,\Theta,\Xi\right)$ is a $C^*$-{\it dynamical system},
meaning that the vector group
acts strongly continuously by automorphisms of the (possibly noncommutative) $C^*$-algebra $\A$.
Let us denote by $\Ai$ the family of elements $f$ such that the mapping
$\Xi\ni X\mapsto\Th_X(f)\in\mathcal A$ is
$C^\infty$. It is a dense $^*$-algebra of $\A$ and also a Fr\'echet algebra with the family of semi-norms
\begin{gather}\label{semicar}
\| f\|_\A^{(k)}:=\sum_{|\alpha|\le
k}\frac{1}{|\alpha|!}\|\partial_X^\alpha\left[\Th_X(f)\right]_{X=0}\|_\A \equiv\sum_{|\alpha|\le
k}\frac{1}{|\alpha|!}\|\delta^\alpha (f)\|_\A ,
\qquad k\in\N .
\end{gather}
To quantize the above structure, one keeps the involution unchanged but introduces on $\Ai$ the product
\begin{gather*}%\label{rodact}
f\,\#\,g:=\pi^{-2n}\int_\Xi\int_\Xi dYdZ\,e^{2i[\![Y,Z]\!]} \Th_Y(f) \Th_Z(g) ,
\end{gather*}
suitably def\/ined by oscillatory integral techniques~\cite{Fol89,Rie1}. One gets a $^*$-algebra $(\Ai,\#,^*)$, which admits a
$C^*$-completion~$\AA$ in a $C^*$-norm $\|\cdot\|_\AA$ def\/ined by Hilbert module techniques~\cite{Rie1}.
The action~$\Th$ leaves~$\Ai$ invariant and extends to a strongly continuous action of the
$C^*$-algebra~$\AA$, that will also be denoted by~$\Th$. The space $\AAi$ of $C^\infty$-vectors coincides with~$\Ai$
and it is a~Fr\'echet space with semi-norms
\begin{gather}\label{semon}
\| f\|_\AA^{(k)}:=\sum_{|\alpha|\le k}\frac{1}{|\alpha|!}\|\partial_X^\alpha\left[\Th_X(f)\right]_{X=0}
\|_\AA \equiv\sum_{|\alpha|\le
k}\frac{1}{|\alpha|!}\|\delta^\alpha (f)\|_\AA ,\qquad k\in\N  .
\end{gather}
By Proposition 4.10 in \cite{Rie1}, there exist $k\in\N $ and $C_k>0$ such that
$\| f\|_\AA \le C_k\| f\|^{(k)}_\A$ for any $f\in\A^\infty=\AA^\infty$.
Replacing here $f$ by $\delta^\alpha f$ for every multi-index $\alpha$, it follows that on $\A^\infty$ the topology given by
the semi-norms~(\ref{semicar}) is f\/iner than the one given by the semi-norms~(\ref{semon}). As a~consequence of
Theorem~7.5 in~\cite{Rie1}, the role of the $C^*$-algebras~$\A$ and $\AA$ can be reversed: one obtains~$\A$ as the deformation
of~$\AA$ by replacing the skew-symmetric form $[\![\cdot,\cdot]\!]$ by $-[\![\cdot,\cdot]\!]$.
Thus~$\A^\infty$ {\it and $\AA^\infty$ coincide as Fr\'echet spaces}.

The deformation transfers to $\Xi$-morphisms. Let
$\left(\A_j,\Theta_j,\Xi,[\![\cdot,\cdot]\!]\right)$, $j=1,2$, be two classical data and
let $\R:\A_1\to\A_2$ be a $\Xi$-morphism, i.e.\ a $C^*$-morphism intertwining the
two actions~$\Theta_1$,~$\Theta_2$. Then $\R$ sends $\A^\infty_1$ into $\A^\infty_2$ and extends to a morphism
$\RR:\AA_1\to\AA_2$ that also intertwines the corresponding actions. In this way, one obtains a
covariant functor.
The functor is exact~\cite{Rie1}: it preserves short exact sequences of~$\Xi$-morphisms. Namely, if~$\J$ is a~(closed,
self-adjoint, two-sided) ideal in $\A$ that is invariant under $\Theta$, then its deformation $\JJ$ can be
identif\/ied with an invariant ideal in~$\AA$ and the quotient~$\AA/\JJ$ is canonically isomorphic to the
deformation  of the quotient~$\A/\J$  under the natural quotient action.

We will refer to {\it the Abelian case} under the following circumstances:
A continuous action $\Theta$ of $\Xi$ by homeomorphisms of the locally compact Hausdorf\/f space~$\Sigma$ is given. For
$(\sigma,X)\in\Sigma\times\Xi$ we are going to use all the notations
$\Theta(\sigma,X)=\Theta_X(\sigma)=\Theta_\sigma(X)\in\Sigma$ for the $X$-transformed of the point~$\sigma$.
The function $\Theta$ is continuous and the homeomorphisms
$\Th_X$, $\Th_Y$ satisfy $\Th_X\circ\Th_Y=\Th_{X+Y}$ for every $X,Y\in\Xi$.

We denote by $\C(\Sigma)$ the Abelian $C^*$-algebra of all complex continuous functions on $\Sigma$
that are arbitrarily small outside large compact subsets of $\Sigma$. When~$\Sigma$ is compact,
$\C(\Si)$ is unital. The action $\Theta$ of $\Xi$ on $\Sigma$ induces an action of~$\Xi$ on
$\C(\Si)$ (also denoted by~$\Theta$) given by~$\Theta_X(f):=f\circ\Theta_X$.
This action is strongly continuous, i.e.\ for any $f\in\C(\Sigma)$ the mapping
\begin{gather}\label{ciuca}
\Xi\ni X\mapsto\Theta_X(f)\in\C(\Sigma)
\end{gather}
is continuous; thus we are placed in the setting presented above.
We denote by $\C(\Si)^\infty\equiv\C^\infty(\Sigma)$ the set of elements $f\in\C(\Sigma)$ such that the
mapping \eqref{ciuca} is $C^\infty$; it is a dense $^*$-algebra of~$\C(\Sigma)$.
The general theory supplies a $C^*$-algebra $\AA\equiv\CC(\Si)$ (noncommutative in general),
acted continuously by the group~$\Xi$, with
smooth vectors $\CC^\infty(\Si)=\C^\infty(\Si)$.


\section[Families of $C^*$-algebras]{Families of $\boldsymbol{C^*}$-algebras}\label{secintra}

Now we give a short review of $\C(T)$-algebras and semi-continuous f\/ields of $C^*$-algebras
(see \cite{Bl,La,Lee,Ni,Rie,Wi} and references therein), outlining the connection between the two notions.

If $\B$ is a $C^*$-algebra, we denote by $\mathcal M(\B)$ its multiplier algebra and by $\mathcal Z\mathcal M(\B)$ its center.
If $\mathcal B_1$, $\mathcal B_2$ are two vector subspaces of $\mathcal M(\B)$, we set $\B_1\cdot\B_2$ for the vector
subspace generated by $\{b_1 b_2\mid b_1\in\B_1,b_2\in\B_2\}$.
We are going to denote by $\C(T)$ the $C^*$-algebra of all complex continuous functions on the (Hausdorf\/f)
locally compact space $T$ that decay at inf\/inity.

\begin{Definition}\label{cedete}
We say that $\B$ is a {\it $\C(T)$-algebra}
if a nondegenerate monomorphism $\Q:\C(T)\rightarrow\mathcal Z\mathcal M(\B)$ is given.
\end{Definition}

We recall that nondegeneracy means that the ideal $\Q[\C(T)]\cdot\B$
is dense in~$\B$.

\begin{Definition}\label{niulaif}
By {\it upper-semi-continuous field of $C^*$-algebras} we mean a family of epimorphisms of $C^*$-algebras
$\big\{\mathcal B\overset{\P(t)}{\longrightarrow}\mathcal B(t)\mid t\in T\big\}$
indexed by a locally compact topological space~$T$ and satisfying:
\begin{enumerate}\itemsep=0pt
\item
For every $b\in\B$ one has $\| b\|_\B=\sup_{t\in T}\| \P(t)b\|_{\B(t)}$.
\item
For every $b\in\B$ the map
$T\ni t\mapsto\| \P(t)b\|_{\B(t)}$ is upper-semi-continuous and decays at inf\/inity.
\item
There is a multiplication $\C(T)\times\B\ni(\varphi,b)\rightarrow\varphi\ast b\in\B$ such that
\[
\P(t)[\varphi\ast b]=\varphi(t) \P(t)b ,\qquad \forall\,t\in T,\quad \varphi\in\C(T) ,\quad  b\in\B .
\]
\end{enumerate}
If, in addition, the map $t\mapsto\| \P(t)b\| $ is continuous for every $b\in \B$, we say that
\[
\big\{\mathcal B\overset{\P(t)}{\longrightarrow}\mathcal B(t)\mid t\in T\big\}
\] is {\it a~continuous field of
$C^*$-algebras}.
\end{Definition}

The requirement~2 is clearly equivalent with the condition that
for every $b\in\B$ and every $\epsilon>0$ the subset $\{t\in T\mid\| \P(t)b\|_{\B(t)}\ge \epsilon\}$  is compact.
One can rephrase~1 as $\cap_t\ker[\P(t)]=\{0\}$,
so one can identify $\B$ with a $C^*$-algebra of sections of the f\/ield; this make the connection with other approaches,
as that of~\cite{Ni} for example. It will always be assumed that $\B(t)\ne\{0\}$ for all $t\in T$.

We are going to describe brief\/ly in which way the two def\/initions above are actually equivalent.

First let us assume that $\B$ is a $\C(T)$-algebra and denote by $\C_t(T)$ the ideal of all the functions in
$\C(T)$ vanishing at the point $t\in T$. We get ideals
$\mathcal I(t):=\overline{\Q\left[\C_t(T)\right]\cdot\B}$ in $\B$, quotients $\B(t):=\B/\mathcal I(t)$ as well as canonical
epimorphisms $\P(t):\B\rightarrow \B(t)$. One also sets
\begin{gather*}%\label{invent}
\varphi\ast b:=\Q(\varphi)b ,\qquad \forall\,\varphi\in\C(T),\quad b\in\B .
\end{gather*}
Then $\big\{\mathcal B\overset{\P(t)}{\longrightarrow}\mathcal B(t)\mid t\in T\big\}$ is an upper-semi-continuous
f\/ield of $C^*$-algebras with multiplication~$\ast$.

Conversely, if an upper-semi-continuous f\/ield
$\big\{\mathcal B\overset{\P(t)}{\longrightarrow}\mathcal B(t)\mid t\in T\big\}$ is given,
also involving the multiplication~$\ast$, we set
\begin{gather*}%\label{invent}
\Q: \ \C(T)\rightarrow\mathcal Z\mathcal M(\B),\qquad \Q(\varphi)b:=\varphi\ast b.
\end{gather*}
In this way one gets a $\C(T)$-algebra and each of the quotients $\B/\mathcal I(t)$ is isomorphic to the f\/i\-ber~$\B(t)$.

To discuss lower semi-continuity we need ${\rm Prim}(\B)$, the space of all the primitive ideals
(kernels of irreducible representations) of $\B$. The hull-kernel topology
turns ${\rm Prim}(\B)$ into a locally compact (not necessarily Hausdorf\/f) topological space. We recall that {\it the hull application}
$\mathcal J\mapsto h(\mathcal J):=\{\mathcal K\in {\rm Prim}(\B)\mid \mathcal J\subset\mathcal K\}$
realizes a containment reversing bijection between the family of ideals of $\B$ and the family
of closed subsets of ${\rm Prim}(\B)$. Its inverse is {\it the kernel map}
$\Omega\mapsto k(\Omega):=\cap_{\mathcal K\in\Omega}\mathcal K$, which is also decreasing.

The Dauns--Hof\/fman theorem establishes the existence of a unique isomorphism
\[
\Gamma: \ BC[{\rm Prim}(\B)]\to \mathcal Z\mathcal M(\B) ,
\] where $BC[{\rm Prim}(\B)]$ is the $C^*$-algebra of bounded and continuous
functions over ${\rm Prim}(\B)$, such that for each $\K\in {\rm Prim}(\B)$,
$\Psi\in BC[{\rm Prim}(\B)]$ and $b\in\B$ we have $\Gamma(\Psi)b+\mathcal K=\Psi(\mathcal K)b+\mathcal K$.
For a~detailed study of the space ${\rm Prim}(\B)$ and a
proof of the Dauns--Hof\/fman theorem, cf.\ Sections~A.2 and~A.3 in~\cite{RW}.
Let us suppose that there is a continuous surjective function $q:{\rm Prim}(\B)\to T$. Then we can def\/ine
$\Q:\C(T)\to \mathcal Z\mathcal M(\B)$ by $\Q(\varphi)=\Gamma(\varphi\circ q)$ and one can check that
$\Q$ endows $\B$ with the structure of a $\C(T)$-algebra.

On the other hand, if we have a nondegenerate monomorphism $\Q:\C(T)\to \mathcal Z\mathcal M(\B)$, we can def\/ine canonically
a continuous map $q:{\rm Prim}(\B)\to T$. One has $q(\mathcal K)=t$ if and only if $\mathcal I(t)\subset \mathcal K$,
and we can recover $\Q$ from the above construction.
Moreover {\it the map $T\ni t\to\| b(t)\|_{\B(t)}\in\mathbb R_+$ is continuous
for every $b\in\B$ $($so we have a continuous field of $C^*$-algebras$)$ if and only if $q$ is open}.
For the proof of this facts see Propositions~C.5 and~C.10 in~\cite{Wi}.



\section[Covariant $\C(T)$-algebras and upper-semi-continuity under Rieffel quantization]{Covariant $\boldsymbol{\C(T)}$-algebras and upper-semi-continuity\\ under Rief\/fel quantization}\label{isra}


Let $T$ be a locally compact Hausdorf\/f space and $(\A,\Th,\Xi,[\![\cdot,\cdot]\!])$ a classical data.
The canonical $C^*$-dynamical system def\/ined by Rief\/fel quantization is $(\AA,\Th,\Xi)$.

\begin{Definition}\label{cedete+}
We say that $\A$ is a {\it covariant $\C(T)$-algebra with respect to the action} $\Th$
if a~nondegenerate monomorphism $\Q:\C(T)\rightarrow\mathcal Z\mathcal M(\A)$ is given
(so it is a $\C(T)$-algebra) and in addition one has
\begin{gather}\label{incep}
\Th_X[\Q(\varphi)f]=\Q(\varphi)\left[\Th_X(f)\right],\qquad \forall\,f\in\A,\quad X\in\Xi,\quad \varphi\in\C(T) .
\end{gather}
\end{Definition}

We intend to prove that the Rief\/fel quantization transforms covariant $\C(T)$-algebras into covariant $\C(T)$-algebras.
For this and for a further result identifying the emerging quotient algebras, we are going to need

\begin{Lemma}\label{densly}
Let $I$ be an ideal of $\C(T)$ and denote by $\overline{\Q(I)\cdot\A}$ the closure of $\Q(I)\cdot\A$ in the $C^*$-algebra~$\A$.
Then $\Q(I)\cdot\A^\infty$ is dense in
$\big(\overline{\Q(I)\cdot\A}\big)^\infty\equiv\big(\overline{\Q(I)\cdot\A}\big)\cap\A^\infty$
for the Fr\'echet topology inherited from~$\A^\infty$.
\end{Lemma}

\begin{proof}
By the covariance condition $\overline{\Q(I)\cdot\A}$ is an invariant ideal of $\A$.

The proof uses regularization. Consider {\it the integrated form of} $\Th$, i.e.\ for each $\Phi\in C_c^\infty(\Xi)$
(compactly supported smooth function) and $g\in\A$ def\/ine
\[
\Th_\Phi(g)=\int_\Xi dY\Phi(Y)\Th_Y(g).
\]
Note that for every $X\in\Xi$ one has
\[
\Th_X\left[\Th_\Phi(g)\right]=\int_\Xi dY\Phi(Y-X)\Th_Y(g).
\]
Then $\Th_\Phi(g)\in\A^\infty$ and for each multi-index $\mu$ we have
\[
\delta^{\mu}\left[\Th_\Phi(g)\right]=(-1)^{|\mu|}\Th_{\partial^{\mu}\Phi}(g)\qquad \text{and}\qquad
\|\delta^{\mu}\left[\Th_\Phi(g)\right]\|_\A \leq \|\partial^{\mu}\Phi\|_{L^1(\Xi)}
\| g\|_\A .
\]
One of the deepest theorems about smooth algebras, the Dixmier--Malliavin theorem~\cite{DM}, say that~$\A^\infty$
is generated (algebraically) by the set of all the elements of the form  $\Th_\Phi(g)$ with $\Phi\in C_c^\infty(\Xi)$ and
$g\in\A$. Replacing~$\A$ with $\overline{\Q(I)\cdot\A}$,
for $f\in\big(\overline{\Q(I)\cdot\A}\big)^\infty$ there exist $\Phi_1,\dots,\Phi_m\in C_c^\infty(\Xi)$ and
$f_1,\dots,f_m\in\overline{\Q(I)\cdot\A}$ such that $f=\sum\limits_{i=1}^m\Th_{\Phi_i}(f_i)$.
Let $\epsilon >0$ and f\/ix a multi-index $\alpha$. Choose $g_1,\dots ,g_m\in\Q(I)\cdot\A$ such that for each $i$
\[
\| f_i-g_i\|_\A \le\frac{\epsilon}{m\|\partial^\alpha\Phi_i\|_{L^1(\Xi)}} .
\]
Then
\begin{gather*}
\left\| \delta^\alpha \left(f-\sum_{i=1}^m\Th_{\Phi_i}(g_i)\right)\right\|_\A=
\left\|\sum_{i=1}^m\Th_{\partial^\alpha\Phi_i}(f_i-g_i)\right\|_\A
\leq
\sum_{i=1}^m\|\partial^\alpha\Phi_i\|_{L^1(\Xi)}\| f_i-g_i\|_\A \le\epsilon .
\end{gather*}
Thus we only need to prove that for each $\Phi\in C_c^\infty(\Xi)$ and $g\in\Q(I)\cdot\A$ the element
$\Th_\Phi(g)$ belongs to $\Q(I)\cdot\A^\infty$. Let $\varphi_1,\dots,\varphi_j\in I$ and
$h_1,\dots,h_j\in\A$ such that $g=\sum\limits_{i=1}^j\Q(\varphi_i)h_i$. Then
\[
\Th_\Phi(g)=\sum_{i=1}^j\Th_\Phi [\Q(\varphi_i)h_i ],
\]
and by covariance, for each index $i$ one has
\begin{gather*}
\Th_\Phi\left[\Q(\varphi_i)h_i\right]=\int_\Xi dY\Phi(Y)\Q(\varphi_i)\Th_X(h_i)=\Q(\varphi_i)
\left[\Th_\Phi(h_i)\right]\in\Q(I)\cdot\A^\infty .\tag*{\qed}
\end{gather*}
\renewcommand{\qed}{}
\end{proof}

\begin{Theorem}\label{crifel}
Rieffel quantization transforms covariant $\C(T)$-algebras into covariant $\C(T)$-algebras: there exists a nondegenerate
monomorphism $\mathfrak Q:\C(T)\rightarrow\mathcal Z\mathcal M(\AA)$ satisfying for all $\varphi\in\C(T)$, $f\in\A$ and
$X\in\Xi$ the covariance relation
$\Th_X[\QQ(\varphi)f]=\QQ(\varphi) [\Th_X(f) ]$.
\end{Theorem}

\begin{proof}
The action $\Th$ of $\Xi$ on $\A$ extends canonically to an action by automorphisms of the multiplier algebra
$\mathcal M(\A)$, also denoted by $\Th$, which is not strongly continuous in general.
But, tautologically, it restricts to a strongly continuous action
$\Th:\Xi\rightarrow{\rm Aut}[\mathcal M_0(\A)]$ on the $C^*$-subalgebra
\begin{gather*}%\label{clar}
\mathcal M_0(\A):=\{m\in\mathcal M(\A)\mid\Xi\ni X\mapsto\Th_X(m)\in\mathcal M(\A)\ {\rm is\ norm\ continuous}\}.
\end{gather*}
In these terms, the covariance condition on $\Q$ says simply that for any $\varphi\in\C(T)$ the multip\-lier~$\Q(\varphi)$ is a f\/ixed point for all the automorphisms~$\Th_X$
(take $f=1$ in~(\ref{incep})). As a very weak consequence one has
$\Q[\C(T)]\subset\mathcal M_0(\A)^\infty$, with an obvious notation for the smooth vectors.

Proposition~5.10 from \cite{Rie1} applied to the unital $C^*$-algebra $\mathcal M_0(\A)$ says that the Rief\/fel quantization of
$\mathcal M_0(\A)$ is a $C^*$-subalgebra of $\mathcal M(\AA)$. Consequently one has $\Q[\C(T)]\subset\mathcal M_0(\A)^\infty
\subset\mathcal M(\AA)$ and this supplies a candidate $\QQ:\C(T)\rightarrow\mathcal M(\AA)$.
This is obviously an injective map and the range is only composed of f\/ixed points, which insures covariance.

Let us set for a moment $\mathcal M:=\mathcal M_0(\A)$, with multiplication~$\cdot$, and denote by
$\mathfrak M\subset\mathcal M(\AA)$ its Rief\/fel quantization, with multiplication legitimately denoted by~$\#$.
For smooth elements $m,n\in\mathcal M^\infty=\mathfrak M^\infty$, one of them being a f\/ixed point central in~$\mathcal M$,
one has $m\#n=m\cdot n=n\cdot m=n\#m$ (Corollary~2.13 in~\cite{Rie1}).
Thus the mapping~$\QQ$ is again a monomorphism and its range is contained in $\mathcal Z\mathfrak M$.
A density argument with respect to the strict topology implies that every $\mathfrak Q(\varphi)$
commutes with all the elements of~$\mathcal M(\AA)$, thus $\QQ[\C(T)]\subset\mathcal Z\mathcal M(\AA)$as required.

Now we only need to show nondegeneracy, i.e.\ the fact that $\QQ[\C(T)]\cdot\AA$ is dense in~$\AA$.
We show the even stronger assertion that $\Q[\C(T)]\cdot\A^\infty=\QQ[\C(T)]\cdot\AA^\infty$ is dense in~$\AA$.
This would follow if we knew that $\Q[\C(T)]\cdot\A^\infty$ is dense in $\AA^\infty$ with respect to its Fr\'echet topology
given by the semi-norms~(\ref{semon});
then we use denseness of $\AA^\infty$ in the weaker $C^*$-norm topology of~$\AA$.

We recall from Section~\ref{sectra} that $\A^\infty$ and $\AA^\infty$ coincide even as Fr\'echet spaces.
Therefore one is reduced to showing that $\Q[\C(T)]\cdot\A^\infty$ is dense in $\A^\infty$ for its Fr\'echet topology.
Taking $\mathcal I=\C(T)$ in Lemma~\ref{densly}, we f\/ind out that $\Q[\C(T)]\cdot\A^\infty$ is dense in
$\big(\overline{\Q[\C(T)]\cdot\A}\big)\cap\A^\infty$, which equals~$\A^\infty$ since $\Q$ has been assumed
nondegenerate. This f\/inishes the proof.
\end{proof}

If $\A$ is a covariant $\C(T)$-algebra, then $\mathcal I(t):=\overline{\Q\left[\C_t(T)\right]\cdot\A}$
is an invariant ideal of~$\A$. We can apply Rief\/fel quantization to
$\mathcal I(t)$, to $\A(t):=\A/\mathcal I(t)$ (with the obvious actions of $\Xi$) and to the projection
$\mathcal P(t):\A\rightarrow\A(t)$. One gets $C^*$-algebras~$\mathfrak I_t$, $\AA_t$ as well as the
morphism $\mathfrak P_t:\AA\rightarrow\AA_t$. By Theorem~7.7 at~\cite{Rie1} the kernel of~$\mathfrak P_t$ is~$\mathfrak I_t$, so~$\AA_t$ can be identif\/ied to the quotient~$\AA/\mathfrak I_t$.

On the other hand, by using the $\C(T)$-structure of the $C^*$-algebra $\AA$ given by Theorem~\ref{crifel}, we have ideals
$\mathfrak I(t):=\overline{\QQ\left[\C_t(T)\right]\cdot\AA}$ in $\AA$ as well as quotients $\AA(t):=\AA/\mathfrak I(t)$
to which we associates projections $\AA\overset{\PP(t)}{\longrightarrow}\AA(t)$.
However, one gets

\begin{Proposition}\label{pieligroasa}
With notation as above, for each $t\in T$ we have $\mathfrak I(t)=\mathfrak I_t$.
In particular, the fibers $\AA(t)=\AA/\mathfrak I(t)$ of  the $\C(T)$-algebra $\AA$ are isomorphic
to the Rieffel quantization $\AA_t$ of the fibers $\A(t)=\A/\mathcal I(t)$ of $\A$ and for each
$f\in\AA$ the mapping $t\mapsto\| \mathfrak P(t)f\|_{\AA(t)}=\| \PP_t f\|_{\AA_t}$
is upper-semi-continuous.
\end{Proposition}

\begin{proof}
We recall that $\mathcal I(t)^\infty$ and $\mathfrak I(t)^\infty$ coincide as Fr\'echet spaces.
By Lemma \ref{densly}, $\QQ\left[\C_t(T)\right]\cdot\AA^\infty$ is dense in $\mathfrak I(t)^\infty$,
thus in $\mathfrak I(t)$, and
$\Q\left[\C_t(T)\right]\cdot\A^\infty$ is dense in $\mathcal I(t)^\infty=\mathfrak I(t)^\infty$,
thus also dense in~$\mathfrak I_t$.

By construction one has $\QQ\left[\C_t(T)\right]\cdot\AA^\infty=\Q\left[\C_t(T)\right]\cdot\A^\infty$;
consequently $\mathfrak I(t)=\mathfrak I_t$ for every $t\in T$ and the proof is f\/inished.
\end{proof}
